\providecommand{\pdfxopt}{a-1b,mathxmp}

%\documentclass[10pt]{article}
%\usepackage{beamerarticle}
\documentclass[ignorenonframetext,slidestop]{beamer}

\input{/home/asreeve/current/class/ers102/pre_latex/luatex_pre.tex}

\setbeamertemplate{enumerate item}[square]
\setbeamertemplate{enumerate subitem}[circle]
\setbeamertemplate{enumerate subsubitem}[square]
\setbeamercolor*{enumerate item}{fg=red,bg=blue}
\setbeamercolor*{enumerate subitem}{fg=blue,bg=red}

\begin{document}

\begin{frame}{Igneous Rocks and Minerals}
  \begin{enumerate}
  \item Igneous Rocks form from a liquid melt (magma).
  \item Originate from fluids deep in the earth or for rocks that have been heated and \textbf{completely} melted
  \item Made up of minerals
    \begin{enumerate}
    \item inorganic
    \item natural
    \item definite chemical composition
    \item crystalline
    \item homogenous
    \end{enumerate}
  \end{enumerate}
\end{frame}

\begin{frame}{Igneous Rock Classification}
  \begin{enumerate}
  \item{Texture}
    \begin{enumerate}
    \item Glass
    \item Aphanitic
    \item Phaneritic
    \item Vesicular (foamy)
    \end{enumerate}
    
  \item{Mineralogy/Chemical Comp}
    \begin{enumerate}
    \item Felsic:Light or Pink
    \item Intermediate: gray or salt \& pepper 
    \item Mafic: Dark
    \end{enumerate}
  \end{enumerate}
\end{frame}

\begin{frame}{Identifying mineral hand samples}
  Mineral identification of large hand samples is done by making simple observations and performing simple tests on the sample.
  \begin{enumerate}
  \item Cleavage: how many and angle between cleavage planes
  \item Color: careful, some minerals come in a range of colors (eg: quartz)
  \item Luster: metallic, glassy, etc.
  \item Hardness: easy it is so scratch (not durability). Hardness of common materials used to 'bracket' sample hardness.
  \item Streak: color of powdered mineral, uses a porcelain plate (h = 6.5-7.5). If a mineral is harder than the streak plate, it can powder the
     streak plate (false response).
  \end{enumerate}
\end{frame}

\begin{frame}{Doing the Lab}{Mineral ID}
  \begin{enumerate}
  \item Minerals:numbers, rocks:letters
  \item Watch the videos or examine hand samples and fill in the attached tables.

  \item More rows on the table than there are samples.
  \item Table of minerals contains minerals not in this lab, they will appear in other labs
  \item Mineral ID tables: one for metallic and one for non-metallic minerals
  \item Once properties defined, ID mineral sample and fill in table
  \end{enumerate}
The following minerals are in this lab: hematite, quartz, hornblende, pyroxene, olivine, biotite,
     muscovite, orthoclase, plagioclase
\end{frame}

\begin{frame}{Doing the Lab}{Rock ID}
  \begin{enumerate}
  \item group rocks by grain size and color
  \item compare with rock ID table and graphics in lab handout
  \item One way to do this is to organize the samples in an array:
    \begin{tabular}{c c c c}
      &light& intermediate & dark \\ \hline \hline
      fine& X & X & X \\
      course& X & X & X \\
      glassy&&& \\
      vesicular &&& \\ \hline
    \end{tabular}
  \end{enumerate}
The following rock samples are present: obsidian, pumice,scoria, rhyolite, andesite, basalt, granite, diorite, and gabbro
  
\end{frame}
\end{document}
